% ===========================================================================
\chapter{The program}\label{ch:program}

All the parts are there. This chapter puts them together into a program,
\code{code.nex}, that you run from the command line: give it a task, and
it works on the sample project. Or give it no task, and it asks for one
task after another, carrying the conversation from each to the next, so
that you can say ``now add a check for that''.

\section{Options}

The program's settings come from the command line. Most have sensible
defaults: the model is \code{claude-opus-5-5} at \code{high} effort, a task
may take at most 40 replies, and a command at most two minutes. Words that
are not options make up the task.

\begin{newfile}{code.nex}
-- code.nex: the coding agent, on the command line.
--
--   nex code.nex "make the failing checks pass"     -- one task, then stop
--   nex code.nex                                    -- type tasks one after another
--   nex code.nex --replay sessions/fix_late_fines.json --yes
--
-- By default the agent works in workspace/, a copy of sample_project/
-- (a small lending library with a bug). It asks before it changes a file
-- or runs a command, unless you give --yes.
--
-- Talking to Claude needs ANTHROPIC_API_KEY (or ANTHROPIC_AUTH_TOKEN) in
-- the environment. --replay plays back a recorded session instead, with
-- no key, no network and no cost.

intern Agent
intern Claude_Model
intern Helpers
intern io/Path

-- The settings of one run, read from the command line.
class Options
create
  from_command_line(args: Array[String]) do
    model_id := "claude-opus-5-5"
    effort := "high"
    max_replies := 40
    command_seconds := 120
    task := ""
    from_folder := ""
    dir := ""
    replay := ""
    record := ""
    let words: Array[String] := []
    from
      let i: Integer := 0
    until
      i >= args.length
    do
      let arg: String := args.get(i)
      if arg = "--yes" then
        yes := true
      elseif arg = "--read-only" then
        read_only := true
      elseif arg = "--fresh" then
        fresh := true
      elseif arg = "--full-output" then
        full_output := true
      elseif arg = "--help" or arg = "-h" then
        help := true
      elseif arg.starts_with("--") then
        if i + 1 >= args.length then
          raise "missing value after " + arg
        end
        let value: String := args.get(i + 1)
        case arg of
          "--model"     then model_id := value
          "--effort"    then effort := value
          "--max-turns" then max_replies := value.to_integer
          "--timeout"   then command_seconds := value.to_integer
          "--from"      then from_folder := value
          "--dir"       then dir := value
          "--replay"    then replay := value
          "--record"    then record := value
          else raise "unknown option " + arg + " (try --help)"
        end
        i := i + 1
      else
        words.add(arg)
      end
      i := i + 1
    end
    task := " ".join(words)
  end

feature
  task: String         -- what to do; "" to ask for tasks one by one
  model_id: String     -- which Claude
  effort: String       -- low, medium, high, xhigh or max
  from_folder: String  -- copy this folder into workspace/ and work there
  dir: String          -- work in this folder itself
  replay: String       -- play back this recorded session
  record: String       -- record the session to this file
  yes: Boolean         -- allow changes without asking
  read_only: Boolean   -- allow no changes at all
  fresh: Boolean       -- start workspace/ again from its template
  full_output: Boolean -- print the output of every tool in full
  help: Boolean
  max_replies: Integer -- the most replies one task may take
  command_seconds: Integer  -- how long a command may run

invariant
  positive_limits: max_replies > 0 and command_seconds > 0
end
\end{newfile}

\code{case} \ldots\ \code{of} picks a branch by value. The invariant catches
a \code{--max-turns 0} as soon as the options are read.

The help text lists them all, and \code{fail} is for mistakes on the
command line: it says what is wrong and stops the program with
\code{exit(1)}, the usual signal to whoever ran it that something failed.

\begin{continued}{code.nex}
function usage_text(): String do
  result := "usage: nex code.nex [options] [task]\n\n"
    + "Without a task, it asks for one task after another (an empty line quits).\n\n"
    + "  --replay FILE     play back a recorded session (no API key needed)\n"
    + "  --record FILE     record the session, for --replay\n"
    + "  --from FOLDER     work in workspace/, a fresh copy of FOLDER (default: sample_project)\n"
    + "  --dir FOLDER      work in FOLDER itself; careful, its files will change\n"
    + "  --fresh           start workspace/ again from a fresh copy\n"
    + "  --yes             let the agent edit files and run commands without asking\n"
    + "  --read-only       never let it edit files or run commands\n"
    + "  --model ID        which Claude (default claude-opus-5-5)\n"
    + "  --effort LEVEL    low, medium, high, xhigh or max (default high)\n"
    + "  --max-turns N     the most replies one task may take (default 40)\n"
    + "  --timeout SECS    how long a command may run (default 120)\n"
    + "  --full-output     print every tool's output in full\n"
end

function fail(message: String) do
  console.print_line(message)
  exit(1)
end
\end{continued}

\section{Choosing the model}

The rest of the file is the program itself, statements at the top level,
run in order. First the options, and two mistakes worth catching before
anything happens:

\begin{continued}{code.nex}
let options: Options := create Options.from_command_line(program.arguments)
if options.help then
  console.print_line(usage_text())
  exit(0)
end
if options.dir /= "" and (options.from_folder /= "" or options.replay /= "" or options.record /= "") then
  fail("--dir can't be used with --from, --replay or --record: they work in a fresh copy")
end
if not Claude_Model.EFFORTS.contains(options.effort) then
  fail("--effort must be one of " + ", ".join(Claude_Model.EFFORTS))
end
\end{continued}

Then the model. With \code{--replay}, it is the recorded session, which
also says what the task was and which folder to copy. Otherwise it is
Claude, and there must be a key; without one, the program explains how to
get one, and how to see the agent at work anyway. With \code{--record},
Claude is wrapped in a recorder.

The variable \code{model} needs a value from the start, so it begins as an
empty script, which every branch below replaces. \code{replayer} keeps the
replay as a \code{Replay_Model}, to ask it at the end whether the session
diverged; \code{model} can't be asked that, being only a \code{Model}.

\begin{continued}{code.nex}
-- The model, and the folder the workspace starts as a copy of.
let model: Model := create Scripted_Model.make([])
let task: String := options.task
let template: String := when options.from_folder = "" then "sample_project" else options.from_folder end
let replayer: ?Replay_Model := nil
if options.replay /= "" then
  if not (create Path.make(user_path(options.replay))).is_file() then
    fail("there is no session file " + options.replay)
  end
  let replay: Replay_Model := create Replay_Model.make(user_path(options.replay))
  replayer := replay
  model := replay
  template := replay.template()
  if task = "" then
    task := replay.task()
  end
else
  if not has_credentials() then
    fail("To talk to Claude, set ANTHROPIC_API_KEY to your API key.\n"
      + "To see the agent at work without one, play back a recorded session:\n\n"
      + "  nex code.nex --replay sessions/fix_late_fines.json --yes")
  end
  model := create Claude_Model.from_environment(options.model_id, options.effort)
  if options.record /= "" then
    if task = "" then
      fail("--record needs a task on the command line")
    end
    model := create Recording_Model.make(model, user_path(options.record), task, template)
  end
end
\end{continued}

\section{The workspace, the permission, the agent}

Normally the agent works in \code{workspace/}, a copy of the sample
project. The copy is made fresh when it doesn't exist yet, when you ask
with \code{--fresh}, and always for a replay or a recording, which must
start from the same files every time. Otherwise the workspace is left as
the last run left it, so that you can carry on. \code{--dir} points the
agent at a folder of your own, in place: it will change your files, so use
it on something under version control.

\begin{continued}{code.nex}
-- The workspace.
let folder: String := user_path("workspace")
if options.dir /= "" then
  folder := user_path(options.dir)
  if not (create Path.make(folder)).is_directory() then
    fail("there is no folder " + options.dir)
  end
else
  let copy_again: Boolean := options.fresh or options.from_folder /= "" or options.replay /= ""
    or options.record /= "" or not (create Path.make(folder)).exists()
  if copy_again then
    if not (create Path.make(user_path(template))).is_directory() then
      fail("there is no folder " + template)
    end
    fresh_workspace(user_path(template), folder)
    console.print_line("(workspace/ is a fresh copy of " + template + "/)")
  else
    console.print_line("(working in workspace/ as it is; --fresh starts it again)")
  end
end
let ws: Workspace := create Workspace.make(folder)
\end{continued}

The agent asks before changing anything, unless told otherwise. That
includes replays: a session file holds commands, and playing back a file
you were sent means running its commands. Asking is the safe default.

\begin{continued}{code.nex}
-- Who allows what.
let permission: Permission := create Ask_First.make()
if options.read_only then
  permission := create Allow_Reading.make()
elseif options.yes then
  permission := create Allow_Everything.make()
end

let tools: Array[Tool] := [create List_Files.make(ws), create Read_File.make(ws),
  create Search.make(ws), create Edit_File.make(ws), create Write_File.make(ws),
  create Run_Command.make(ws, options.command_seconds)]
let agent: Agent := create Agent.make(model, tools, permission,
  create Reporter.make(false, options.full_output))
agent.set_max_replies(options.max_replies)
console.print_line("(the model is " + model.name() + ")")
\end{continued}

\section{Running tasks}

A task can fail for reasons that are nobody's fault: the network goes
down, or the API refuses a key. That should end the task, with a message,
not the program: in interactive use you may want to try again.
\code{attempt} catches the failure, and reports it.

\begin{continued}{code.nex}
-- Run one task, reporting any failure of the model or the network
-- instead of stopping.
function attempt(agent: Agent, task: String) do
  do
    agent.run(task)
  rescue
    console.print_line("\nerror: " + exception)
  end
end

if task /= "" then
  attempt(agent, task)
else
  console.print_line("Type a task, or an empty line to quit.")
  let going: Boolean := true
  from until not going do
    let line: ?String := console.read_line("\nyou> ")
    if ?line as text then
      if text.trim() = "" then
        going := false
      else
        attempt(agent, text.trim())
      end
    else
      going := false
    end
  end
end

if ?replayer as replay then
  if replay.diverged_at > 0 then
    console.print_line("\nNote: from reply " + replay.diverged_at.to_string
      + " on, the tools gave different results than when the session was recorded,"
      + " so the recorded replies may not fit what happened this time.")
  end
end
\end{continued}

\section{Running it}

Here is the help:

\runoutput{code.nex --help}{help}

Without a key, the program says what to do:

\runoutput{code.nex "Fix the checks."}{no_key_fails}

So let's play back the session that comes with the tutorial, letting the
agent do what it likes, with \code{--yes}:

\runoutput{code.nex --replay sessions/fix_late_fines.json --yes}{replay}

The same session without \code{--yes} asks before the first command, the
edit and the second command. Here the answers were typed in advance: yes
to the command, ``always'' to the edit, and yes to the second command.
(The answers don't show, because they were not typed at a keyboard.)

% input: y\na\ny
\runoutput{code.nex --replay sessions/fix_late_fines.json}{replay_asking}

After either run, the fixed project is in \code{workspace/}, and
\code{cd workspace; nex checks.nex} shows all seven checks passing.

\section{For real}

With \code{ANTHROPIC_API_KEY} set, the same program talks to Claude:

\begin{shell}
nex code.nex "Some of the checks fail. Find out why and fix it."
nex code.nex --fresh                    # a fresh copy, then type tasks
nex code.nex --read-only "Explain how fines are worked out."
nex code.nex --effort low --model claude-sonnet-5-5 "Add a check for a loan returned early."
nex code.nex --record sessions/mine.json --yes "Some of the checks fail. Find out why and fix it."
nex code.nex --replay sessions/mine.json --yes
\end{shell}

Each real run is different: the model may look at the files in another
order, or explain the bug in other words, or (now and then) take a wrong
turn and come back from it. That is exactly why recorded sessions are
useful.
